\documentclass[12pt,a4paper,oneside]{book}

% =====================================================================
% PACOTES ESSENCIAIS
% =====================================================================
\usepackage[brazilian]{babel}
\usepackage[T1]{fontenc}
\usepackage[utf8]{inputenc}
\usepackage{amsmath}
\usepackage{amssymb}
\usepackage{amsthm}
\usepackage{geometry}
\usepackage{fancyhdr}
\usepackage{hyperref}
\usepackage{booktabs}
\usepackage{mathtools}

% =====================================================================
% GEOMETRIA
% =====================================================================
\geometry{
    left=2.5cm,
    right=2.5cm,
    top=3cm,
    bottom=3cm,
    headheight=14pt
}

% =====================================================================
% CABECALHO E RODAPE
% =====================================================================
\pagestyle{fancy}
\fancyhf{}
\fancyhead[L]{\textit{Capitulo 4: Expansao de Taylor}}
\fancyhead[R]{\thepage}
\fancyfoot[C]{\small Principios de Modelagem Matematica}
\renewcommand{\headrulewidth}{0.4pt}
\renewcommand{\footrulewidth}{0.4pt}

% =====================================================================
% AMBIENTES
% =====================================================================
\theoremstyle{definition}
\newtheorem{definition}{Definicao}[chapter]
\newtheorem{example}[definition]{Exemplo}
\newtheorem{remark}[definition]{Observacao}
\newtheorem{problem}{Problema}

% =====================================================================
% DOCUMENTO
% =====================================================================
\title{\textbf{Capitulo 4: Expansao de Taylor}\\[0.4cm]
\large Principios de Modelagem Matematica}
\author{Baseado em Dym, C. L.}
\date{2026-04-14}

\begin{document}

\maketitle
\tableofcontents
\newpage

\chapter{Expansao de Taylor}

\section{Ideia central do capitulo}

A expansao de Taylor e uma ferramenta para aproximar funcoes complicadas por polinomios em torno de um ponto.
Em modelagem matematica, isso permite:

\begin{itemize}
    \item simplificar modelos nao lineares;
    \item estimar erros de aproximacao;
    \item linearizar sistemas perto de um estado de equilibrio;
    \item obter formulas praticas para calculo e simulacao.
\end{itemize}

\section{Serie de Taylor e Serie de Maclaurin}

\begin{definition}[Serie de Taylor]
Se $f$ tem derivadas ate ordem $n+1$ em torno de $a$, entao
\begin{equation}
f(x)=\sum_{k=0}^{n}\frac{f^{(k)}(a)}{k!}(x-a)^k + R_{n+1}(x),
\end{equation}
onde $R_{n+1}(x)$ e o resto (erro de truncamento).
\end{definition}

\begin{definition}[Serie de Maclaurin]
E o caso particular com $a=0$:
\begin{equation}
f(x)=\sum_{k=0}^{n}\frac{f^{(k)}(0)}{k!}x^k + R_{n+1}(x).
\end{equation}
\end{definition}

\section{Resto de Lagrange e controle de erro}

Uma forma comum do erro e:
\begin{equation}
R_{n+1}(x)=\frac{f^{(n+1)}(\xi)}{(n+1)!}(x-a)^{n+1},
\end{equation}
para algum $\xi$ entre $a$ e $x$.

\begin{remark}
Quanto menor $|x-a|$ e maior a ordem $n$, menor tende a ser o erro.
\end{remark}

\section{Exemplos fundamentais}

\subsection{Exemplo 1: exponencial}

Para $f(x)=e^x$ em torno de $0$:
\begin{equation}
e^x = 1 + x + \frac{x^2}{2!} + \frac{x^3}{3!} + \cdots
\end{equation}

Aproximacao de segunda ordem:
\begin{equation}
e^x \approx 1 + x + \frac{x^2}{2}.
\end{equation}

Para $x=0.1$:
\begin{equation}
e^{0.1} \approx 1 + 0.1 + 0.005 = 1.105,
\end{equation}
enquanto o valor real e aproximadamente $1.105170\ldots$.

\subsection{Exemplo 2: seno}

Para $f(x)=\sin x$ em torno de $0$:
\begin{equation}
\sin x = x - \frac{x^3}{3!} + \frac{x^5}{5!} - \cdots
\end{equation}

Para angulos pequenos, usamos a aproximacao linear:
\begin{equation}
\sin x \approx x.
\end{equation}

Essa aproximacao aparece em pendulo simples e vibracoes de pequena amplitude.

\subsection{Exemplo 3: logaritmo}

Para $|x|<1$:
\begin{equation}
\ln(1+x)=x-\frac{x^2}{2}+\frac{x^3}{3}-\frac{x^4}{4}+\cdots
\end{equation}

Com $x=0.2$ e truncando em ordem 2:
\begin{equation}
\ln(1.2) \approx 0.2 - 0.02 = 0.18.
\end{equation}
Valor real: $\ln(1.2)\approx 0.1823$.

\section{Tabela de series de Taylor mais usadas}

A tabela abaixo reune as series mais frequentes em modelagem matematica.
Todas sao expansoes em torno de $x=0$ (Maclaurin).

\begin{center}
\begin{tabular}{lll}
\toprule
\textbf{Funcao} & \textbf{Expansao ate ordem 3 ou 4} & \textbf{Intervalo de validade} \\
\midrule
$e^x$           & $1 + x + \dfrac{x^2}{2!} + \dfrac{x^3}{3!} + \cdots$ & $x \in \mathbb{R}$ \\[6pt]
$\sin x$        & $x - \dfrac{x^3}{3!} + \dfrac{x^5}{5!} - \cdots$ & $x \in \mathbb{R}$ \\[6pt]
$\cos x$        & $1 - \dfrac{x^2}{2!} + \dfrac{x^4}{4!} - \cdots$ & $x \in \mathbb{R}$ \\[6pt]
$\ln(1+x)$      & $x - \dfrac{x^2}{2} + \dfrac{x^3}{3} - \cdots$ & $|x| < 1$ \\[6pt]
$(1+x)^n$       & $1 + nx + \dfrac{n(n-1)}{2!}x^2 + \cdots$ & $|x| < 1$, $n \in \mathbb{R}$ \\[6pt]
$\tanh x$       & $x - \dfrac{x^3}{3} + \dfrac{2x^5}{15} - \cdots$ & $x \in \mathbb{R}$ \\[6pt]
$(1+x)^{-1}$    & $1 - x + x^2 - x^3 + \cdots$ & $|x| < 1$ \\[6pt]
$\sqrt{1+x}$    & $1 + \dfrac{x}{2} - \dfrac{x^2}{8} + \dfrac{x^3}{16} - \cdots$ & $|x| < 1$ \\[6pt]
\bottomrule
\end{tabular}
\end{center}

\subsection{Aproximacao binomial}

A serie $(1+x)^n$ para qualquer $n$ real com $|x| \ll 1$ reduz-se a:
\begin{equation}
(1+x)^n \approx 1 + nx.
\end{equation}

Casos praticos frequentes:
\begin{itemize}
    \item $\sqrt{1+x} = (1+x)^{1/2} \approx 1 + \dfrac{x}{2}$ \quad (variacao percentual pequena)
    \item $(1+x)^{-1} \approx 1 - x$ \quad (denominador quase constante)
    \item $(1+x)^{-3/2} \approx 1 - \dfrac{3}{2}x$ \quad (campo gravitacional, analise de erros)
\end{itemize}

\begin{remark}
A aproximacao binomial e especialmente util em propagacao de incerteza:
se $f = a^n$ e $a = a_0(1+\epsilon)$ com $\epsilon \ll 1$, entao $f \approx a_0^n(1+n\epsilon)$.
O erro relativo em $f$ e $n$ vezes o erro relativo em $a$.
\end{remark}

\subsection{Pendulo simples: linearizacao via Taylor}

A equacao exata do pendulo de comprimento $\ell$ e:
\begin{equation}
\ddot{\theta} + \frac{g}{\ell}\sin\theta = 0.
\end{equation}

Expandindo $\sin\theta$ em torno de $\theta = 0$:
\begin{equation}
\sin\theta = \theta - \frac{\theta^3}{6} + \cdots \approx \theta \quad \text{(para } |\theta| \ll 1\text{)}.
\end{equation}

O modelo linearizado e:
\begin{equation}
\ddot{\theta} + \frac{g}{\ell}\theta = 0,
\end{equation}
com solucao $\theta(t) = A\cos(\omega t + \phi)$, onde $\omega = \sqrt{g/\ell}$ e periodo:
\begin{equation}
T = \frac{2\pi}{\omega} = 2\pi\sqrt{\frac{\ell}{g}}.
\end{equation}

A analise dimensional (Capitulo 2) ja previa $T \propto \sqrt{\ell/g}$; a serie de Taylor fornece
o coeficiente $2\pi$ e confirma a hipotese de angulos pequenos.

\section{Linearizacao de modelos nao lineares}

Considere uma EDO nao linear:
\begin{equation}
\dot{x}=f(x).
\end{equation}

Perto do ponto de equilibrio $x^*$ com $f(x^*)=0$, definindo perturbacao $\eta=x-x^*$:
\begin{equation}
\dot{\eta} \approx f'(x^*)\eta.
\end{equation}

\begin{remark}
Esse passo transforma um problema nao linear em um problema linear local, facilitando analise de estabilidade.
\end{remark}

\subsection{Exemplo 4: modelo logistico}

Modelo:
\begin{equation}
\dot{P}=rP\left(1-\frac{P}{K}\right).
\end{equation}

Equilibrios: $P^*=0$ e $P^*=K$.

Derivada:
\begin{equation}
f'(P)=r\left(1-2\frac{P}{K}\right).
\end{equation}

No equilibrio $P^*=K$:
\begin{equation}
f'(K)=r(1-2)=-r<0,
\end{equation}
logo, equilibrio estavel localmente.

\section{Aproximacoes uteis em modelagem}

\subsection{Exemplo 5: energia potencial elastica}

Suponha energia potencial nao linear simplificada:
\begin{equation}
U(x)=U_0 + ax + bx^2 + cx^3.
\end{equation}

Perto de $x=0$, o termo dominante (para deslocamentos pequenos) e ate segunda ordem:
\begin{equation}
U(x) \approx U_0 + ax + bx^2.
\end{equation}

Se $a=0$ (equilibrio no ponto), obtemos comportamento tipo mola linear.

\subsection{Exemplo 6: aproximacao de raiz}

Para $f(x)=\sqrt{1+x}$ em torno de $0$:
\begin{equation}
\sqrt{1+x}=1+\frac{x}{2}-\frac{x^2}{8}+\frac{x^3}{16}-\cdots
\end{equation}

Aproximacao de primeira ordem:
\begin{equation}
\sqrt{1+x}\approx 1+\frac{x}{2}.
\end{equation}

Muito usada para variacoes pequenas em engenharia.

\section{Roteiro pratico para aplicar Taylor}

\begin{enumerate}
    \item Escolha o ponto de expansao (geralmente equilibrio ou ponto de interesse).
    \item Calcule derivadas necessarias.
    \item Defina a ordem de truncamento conforme precisao desejada.
    \item Estime o resto para controlar erro.
    \item Verifique intervalo de validade (proximidade do ponto de expansao).
\end{enumerate}

\section{Exercicios curtos}

\begin{problem}
Use Maclaurin de ordem 3 para aproximar $e^{-0.2}$.
\end{problem}

\begin{problem}
Linearize $f(x)=\cos x$ em torno de $x_0=\pi/4$.
\end{problem}

\begin{problem}
Aproxime $\ln(1+x)$ em ordem 2 e estime o erro para $x=0.1$ usando o resto de Lagrange.
\end{problem}

% =====================================================================
% EXEMPLO TIPO PROVA
% =====================================================================
\section{Exemplo tipo prova: aproximacao binomial e taxa de decaimento}

\begin{problem}[Prova -- IIo semestre 2025]
O valor de um carro novo e $V_0 = R\$\,51{.}500$ e o veiculo perde $12\%$ de seu valor
por ano, de forma que $V(t) = V_0\,(0{,}88)^t$.
\begin{enumerate}
    \item Use a aproximacao binomial (Taylor ordem 1) para estimar a taxa de
    depreciacao mensal.
    \item Calcule o erro relativo da aproximacao em relacao ao valor exato.
    \item Escreva a equacao diferencial equivalente e determine a constante de decaimento
    $\lambda$ (em anos$^{-1}$).
\end{enumerate}
\end{problem}

\textbf{Solucao}:

\textbf{1. Aproximacao binomial}:

A taxa mensal e obtida de $(0{,}88)^{1/12} = (1 - 0{,}12)^{1/12}$.
Usando a aproximacao binomial $(1+x)^n \approx 1 + nx$ com $x = -0{,}12$ e $n = 1/12$:
\begin{equation}
(1 - 0{,}12)^{1/12} \approx 1 + \frac{1}{12}(-0{,}12) = 1 - 0{,}01 = 0{,}99.
\end{equation}
Taxa de depreciacao mensal $\approx 1\%$ ao mes.

\textbf{2. Valor exato e erro relativo}:
\begin{equation}
(0{,}88)^{1/12} = e^{\ln(0{,}88)/12} \approx e^{-0{,}12783/12} = e^{-0{,}01065} \approx 0{,}9894.
\end{equation}
Taxa exata: $1 - 0{,}9894 = 1{,}06\%$ ao mes.
Erro relativo: $\dfrac{|1{,}06\% - 1\%|}{1{,}06\%} \approx 5{,}7\%$.

\textbf{3. Equacao diferencial e constante $\lambda$}:

Reescrevendo $V(t) = V_0\,e^{-\lambda t}$, identificamos $\lambda$ por:
\begin{equation}
e^{-\lambda} = 0{,}88 \implies \lambda = -\ln(0{,}88) = \ln\!\left(\frac{1}{0{,}88}\right) \approx 0{,}128\;\text{ano}^{-1}.
\end{equation}
A equacao governante e $\dot{V} = -\lambda V$, com solucao $V(t) = V_0\,e^{-0{,}128\,t}$.

\begin{remark}
A aproximacao binomial fornece $1\%$ ao mes, enquanto o valor exato e $1{,}06\%$ ao mes
--- erro de $5{,}7\%$. Para maior precisao, usa-se a segunda ordem:
$(1+x)^n \approx 1 + nx + \frac{n(n-1)}{2}x^2$, que daria
$1 - 0{,}01 - \frac{1}{288}(0{,}12)^2 \approx 1 - 0{,}0106$, ou seja, $1{,}06\%$. $\checkmark$
\end{remark}

\section{Resumo final}

A expansao de Taylor e um pilar da modelagem porque permite substituir funcoes complicadas por expressoes polinomiais controladas por erro.
Isso facilita analise, simulacao e interpretacao fisica de sistemas reais.

Ideias-chave do capitulo 4:
\begin{itemize}
    \item aproximacao local de funcoes;
    \item erro de truncamento e validade da aproximacao;
    \item linearizacao de sistemas nao lineares;
    \item uso pratico em problemas de engenharia e fisica.
\end{itemize}

\end{document}
